\subsection{Teilaufgabe 2}
\subsubsection{Aufgabenstellung}
Wir sollen ein Programm schreiben welches welches 2 nxn Matrizen miteinander Addieren und Multiplizieren kann.

\subsubsection{Anforderungsdefinition}
\begin{enumerate}
	\item Addiere zwei nxn Matrizen.
	\item Multipliziere zwei nxn Matrizen.
\end{enumerate}

\subsubsection{Entwurf}
\includegraphics[scale=0.75]{uml/uml_c4_p2.png}

\subsubsection{Quellcode}
\paragraph{Matrizen.java}\
\lstinputlisting[language = Java , frame = trBL , escapeinside={(*@}{@*)}]{../chapter_04/src/chapter_04/Matrizen.java}

\subsubsection{Testdokumentation}
Das Programm hat nach dem Aufruf Zwei 2x2 Matrizen erstellt und initialisiert, anschließend 
miteinander Addiert und Multipliziert. Dabei kam das Richtige Ergebnis raus.

\subsubsection{Benutzungshinweise}
Navigieren Sie in der Kommandozeile zum dem Ordner, wo sich die Java Datei befindet.
Danach führen sie "javac Matrizen.java\dq \space auf. Jetzt können Sie das Programm mit
"java Matrizen\dq \space starten. Nach dem das Programm gestartet ist, können Sie die
grö\ss e der Matrix angeben.

\subsubsection{Anwendungsbeispiel}
Nach dem Aufruf von java Matrizen, sollten wir nun folgendes sehen:
\begin{lstlisting}[frame = trBL , escapeinside={(*@}{@*)}]
[sebastian@laptop bin]$ java Matrizen
Matrix A:
70	50	
16	52	

Matrix B:
80	75	
11	33	

Addition von A und B:
150	125	
27	85	

Multiplikation von A und B:
6150	6900	
1852	2916
[sebastian@laptop bin]$  
\end{lstlisting}
